\section{Introduction and the connection between the proofs}\label{sec:intro}
The continuum limit of a lattice Gaussian field is conceptually elementary:
its finite-dimensional laws are Gaussian and its covariance has a Fourier
symbol. A proof on a space of distributions nevertheless needs more than
pointwise convergence of that symbol. It must place every lattice law on a
common measurable carrier, prove tightness there, and show that Fourier data
determine every possible limit. These requirements are especially visible
when the embedding is a sum of Dirac masses rather than an interpolation.

This article compares two completed instances of that problem. In two
spatial dimensions the formal construction controls an $H^{-2}$ energy and
proves weak convergence on a concrete $H^{-3}$ carrier~\cite{aiXiv2D}. In four
dimensions
it controls an $H^{-3}$ energy and proves convergence on an $H^{-4}$ carrier
to an explicitly constructed Gaussian-series law. These are related proofs,
not a claim that changing a dimension parameter in the first Lean theorem
produces the second. The four-dimensional summability, carrier maps and
real-coordinate reconstruction have their own proofs.

Our main mathematical synthesis is the following statement. We use a Fourier
Sobolev realization defined precisely in Section~\ref{sec:carrier}.
\begin{theorem}[Common paper theorem]\label{thm:main}
Let $d\geq1$ be an integer, $L>0$, and $m>0$. Put $a=L/N$ for positive
integers $N$. On $(\Z/N\Z)^d$ let $\phi_N$ be the centered real Gaussian with
site covariance $a^{-d}(m^2-\Delta_a)^{-1}$, and define
\begin{equation}\label{eq:lift}
 X_N=a^d\sum_x\phi_N(x)\delta_{[ax]}.
\end{equation}
For every $s>d/2$, the laws of $X_N$ converge weakly on $\HS{s}(\T_L^d)$
to the unique centered real periodic Gaussian field with covariance
\begin{equation}\label{eq:target-cov}
 Q(f,g)=L^{-d}\sum_{k\in\Z^d}
 \frac{\overline{Ff(k)}Fg(k)}{m^2+(2\pi/L)^2\abs{k}_2^2},\qquad
 Ff(k)=\int_{\T_L^d}f(x)e^{-2\pi i k\cdot x/L}\dd x
\end{equation}
for real smooth periodic $f,g$.
\end{theorem}
\evidence{Theorem~\ref{thm:main} is a paper theorem proved below for arbitrary
$d$ and admissible $s$. The public Lean library verifies the fixed choices
$(d,r,s)=(2,2,3)$ and $(4,3,4)$, where $r$ denotes the stronger moment space.
It does not contain an arbitrary-dimensional formalization of this theorem.}

The useful connection is the chain
\begin{equation}\label{eq:architecture}
\begin{gathered}
\text{finite lattice covariance and raw Fourier coefficients}\\
\Downarrow\\
\text{uniform stronger-space moment and compact weaker-space balls}\\
\Downarrow\\
\text{tightness, finite Gaussian marginal limits, and cylinder uniqueness}\\
\Downarrow\\
\text{weak convergence of the full sequence}.
\end{gathered}
\end{equation}
Dimension enters the summability threshold and the compact carrier estimates.
The final probability argument has the same structure in both cases.

\subsection{Context and scope}
Fourier and nearest-neighbour descriptions of Gaussian free fields are
classical; Sheffield's introduction provides background~\cite{Sheffield}.
Discrete-to-continuous limits on tori in arbitrary dimension also occur for
grounded polyharmonic fields and associated random measures~\cite{DSHKS}.
Those models have different covariance powers and grounding conventions.
Our subject is the positive-mass resolvent and the unmodified atomic lift.
We make no priority claim for its classical convergence mechanism.
The contribution here is a comparison of two concrete formal proof chains,
a self-contained normalization and tightness argument, and reproducible
evidence that identifies precisely which statements the Lean kernel checks.

All volumes and masses in the article remain fixed and positive. Neither an
infinite-volume limit nor an interacting four-dimensional field follows from
this result. An interacting two-dimensional construction requires additional
renormalization and uniform integrability arguments. Osterwalder--Schrader
reconstruction is outside the present claims.
